\documentclass[11pt]{article}

\usepackage{amsmath,amssymb,amstext,amsthm}
\usepackage{geometry}
\usepackage{fancyhdr}
\usepackage[pdftex]{graphicx}
\usepackage{xcolor}

\geometry{
  verbose,
  dvips,
  width=430pt, marginparsep=0pt, marginparwidth=0pt,
  top=90pt, headheight=12pt, headsep=20pt, footskip=30pt, bottom=80pt}

\setlength{\parskip}{\medskipamount}
\setlength{\parindent}{0pt}

\linespread{1.05}

\newtheorem{theorem}{Theorem}
\newtheorem{lemma}[theorem]{Lemma}
\newtheorem{proposition}[theorem]{Proposition}
\newtheorem{corollary}[theorem]{Corollary}
\newtheorem{conjecture}[theorem]{Conjecture}
\newtheorem{obs}{Observation}
\newtheorem{clm}{Claim}
\newtheorem{ques}{Question}
\newtheorem{problem}{Problem}

\newtheorem{ex}{Example}


\newcommand{\spc}{\hspace{0.08in}}
\newcommand{\vs}{\vspace*{.1in}}
\newcommand{\E}{\mathbb{E}}
\newcommand{\Prob}{\mathbb{P}}
\newcommand{\edit}[1]{\emph{\textcolor{blue}{#1}}}


\renewenvironment{proof}{{\noindent \textbf{\textit{Proof.}}}}
{\hfill $\Box$ \vspace*{0.1in}}
\newenvironment{proofc}{{\noindent \textit{Proof of Claim.}}}
{\hfill $\Box$}


\begin{document}

\begin{LARGE}\begin{center}\bf 14.8 Lagrange Multipliers\end{center}
\end{LARGE}

\vs\vs



\section{Warm Up}



\section{Motivation}
What if we do not have a bounded region in which we want to find absolute extrema but rather a specific curve or constraint we must live on? Yes we could try to substitute this curve into our surface and turn the problem into a calc-1 problem, but if the curve is complicated, or implicitly defined, this could get real tricky.


\section{Definitions \& Theorems}

\begin{enumerate}
\item For a \emph{Lagrange multiplier problem} we are always given a surface and a constraint. More specifically our surface may be $z=f(x,y)$ and our constraint may be $g(x,y) = c$. Notice that the constraint is a single level curve while the surface is an infinite family of level curves.

\item The astute observation is the following: At an absolute extrema, the normal of the constraint level curve will be going in the same direction as the normal of a level curve of our original surface. All of this is to say 

\[f_x = \lambda g_x \text{ and } f_y = \lambda g_y\]


We call $\lambda$ the multiplier constant.


\item Note that this idea scales for higher dimensional surfaces!




\end{enumerate}





\newpage

\section{Examples}

\begin{problem}
Use Lagrange multipliers to find the extrema of the below surface given the constraint (you do not need to find the extrema's values).

\vspace{.25cm}

Maximize: $f(x,y) = x^2 + 4y^2 - 2x +8y$

\vspace{.25cm}

Subject to the constraint: $x + 2y = 7$
\end{problem}


\newpage

\begin{problem} Use Lagrange multipliers to find the absolute extrema of the function $f(x,y) = 2x^2 + y^2+x-3y+5$ constrained by the level curve $2x-2y=\frac{11}{2}$.
\end{problem}


\newpage

\section{Scratch Work}

\end{document} 